% ID: SEL-D0071-OB-E01
% ESTADO: cotejada. Segunda pasada visual y borrador compilado revisados.
% Fuente: Selectividad/sel_2009_matematicas_aplicadas/mat apli CC SS05.pdf
% SHA-256: ca51ab435c7584a8e047e19d0826f3273f3e88aed6110ccd00e46a40e85518b1
% Página: 2; opción B; ejercicio original 1.

Una tienda dispone de latas de conserva de tomate de tres fabricantes: A, B y C. El fabricante A envasa el tomate en latas de 250 g, el fabricante B lo envasa en latas de 500 g y el fabricante C en latas de 1 kg. Esas latas de tomate se venden a 1, $1.8$ y $3.3$ euros, respectivamente. Compramos en total 20 latas, que pesan un total de 10 kg y nos cuestan $35.6$ euros. Queremos saber cuántas latas de cada fabricante hemos comprado.
\begin{enumerate}[label=\alph*)]
\item \textbf{(1 punto)} Plantee el sistema de ecuaciones que resolvería el problema anterior.
\item \textbf{(2 puntos)} Resuelva el problema.
\end{enumerate}
